% ============================================================
% ABSTRACT
% ============================================================
Pooled calibration of conformal prediction intervals, as in the
Ensemble Batch Prediction Intervals (EnbPI) algorithm of
\citet{xu2023conformal}, carries a persistent bias under seasonal
heteroskedasticity: the per-season coverage gap contains a term
$2b_{s^*}$, which measures the discrepancy between the season-specific
and the pooled residual distribution and does not vanish as the sample
grows, so that volatile seasons are systematically undercovered.  Two
schemes remove it, and they are the endpoints of a one-parameter
family: stratifying the calibration buffer by season, and standardising
residuals by a season-specific scale before pooling them.  The second
is the more efficient of the two, and it is valid only when the seasons
differ in scale alone; when they differ in the shape of the error
distribution it leaves the coverage spread where pooling left it, while
stratification cuts it by two thirds.

Stratification is usually thought too expensive at the per-season
sample sizes a monthly series provides, and we show that most of that
expense belongs to the interval construction rather than to the
conditioning.  The empirical quantile of a short buffer is biased
inwards, and the minimum-width line search selects the narrowest of
many candidate quantile pairs computed from the same residuals, which
alone costs four to six percentage points of coverage at these buffer
lengths.  Building the interval from conformal order statistics with a
fixed asymmetry repairs both, and stratified coverage at twenty
residuals per season then rises from $0.79$ to its nominal $0.90$ in
four different seasonal designs, the residual cost being width rather
than coverage.  On the food-at-home component of the Brazilian consumer
price index (IPCA; BCB series 1635, $T_s = 20$) the same change takes
coverage from $79.2\%$ to exactly $90.0\%$ with no dispersion across
the twelve calendar months, and on monthly export growth (BCB series
1402, $T_s = 35$) it raises the worst month from $70\%$ to $90\%$ and
overall coverage from $86.7\%$ to $95.0\%$, the overshoot being the
conservatism of an interval whose endpoints are still the extremes of
its buffer.  A permutation test rejects the pure-scale restriction on
both series.

\bigskip

\noindent\textbf{Resumo:}
\begin{otherlanguage}{brazilian}%
A calibração conjunta (\emph{pooled}) de intervalos conformes, como no
algoritmo EnbPI de \citet{xu2023conformal}, carrega um viés persistente
sob heteroscedasticidade sazonal: a lacuna de cobertura por estação
contém um termo $2b_{s^*}$ que não desaparece à medida que a amostra
cresce, de forma que as estações voláteis ficam sistematicamente
subcobertas.  Dois esquemas o eliminam, e são os extremos de uma
família de um parâmetro: estratificar o buffer de calibração por
estação, ou padronizar os resíduos por uma escala específica de cada
estação antes de agrupá-los.  O segundo é o mais eficiente dos dois,
porém só é válido quando as estações diferem apenas em escala; quando
diferem na forma da distribuição dos erros, ele deixa a dispersão da
cobertura onde a calibração conjunta a deixou, ao passo que a
estratificação a reduz em dois terços.

Costuma-se considerar a estratificação cara demais para os tamanhos
amostrais por estação de uma série mensal, e mostramos que a maior
parte desse custo pertence à construção do intervalo, e não ao
condicionamento.  O quantil empírico de um buffer curto é enviesado
para dentro, e a busca por largura mínima seleciona o par de quantis
mais estreito entre muitos candidatos calculados dos mesmos resíduos, o
que sozinho custa de quatro a seis pontos percentuais de cobertura
nesses tamanhos.  Construir o intervalo a partir de estatísticas de
ordem conformes, com assimetria fixa, corrige as duas coisas, e a
cobertura estratificada com vinte resíduos por estação sobe então de
$0{,}79$ para o nominal $0{,}90$ em quatro desenhos sazonais
diferentes, sendo o custo remanescente de largura, e não de cobertura.
No componente de Alimentação no domicílio do IPCA (BCB série 1635,
$T_s = 20$), a mesma mudança leva a cobertura de $79{,}2\%$ para
exatamente $90{,}0\%$, sem dispersão entre os doze meses do calendário;
no crescimento mensal das exportações (BCB série 1402, $T_s = 35$), ela
eleva o pior mês de $70\%$ para $90\%$.  Um teste de permutação rejeita
a restrição de escala pura nas duas séries.
\end{otherlanguage}
